\documentclass{article}
\usepackage{graphicx} % Required for inserting images
\usepackage[utf8]{inputenc} % Per caratteri accentati
\usepackage[T1]{fontenc}    % Per una corretta codifica dei font
\usepackage[italian]{babel} % Per sillabazione in italiano
\usepackage{float}
\usepackage{hyperref}
\usepackage{listings}
\usepackage{color}
\usepackage{amssymb}
\usepackage{amsmath}
\usepackage{multicol}
\usepackage{etaremune}

%VALORI SISTEMABILI
\usepackage[a4paper, top= 1 cm, bottom= 1.5 cm, left= 2 cm, right= 2 cm]{geometry}

\title{Reti di Calcolatori}
\author{Francesco Lanzarini}
\date{ }

\begin{document}

\maketitle

\section{Introduzione}
Per \textbf{rete di calcolatori} si intende un insieme di dispositivi autonomi e interconnessi tra di loro, essa è necessaria per supportare la comunicazione tra utenti e, soprattutto, tra macchine
(\textit{internet of things}). Le principlai comnunicazioni riguardano la condivisione di informazioni, di dispositivi, di risorse, l'accesso a calcolatori remoti e il calcolo distribuito.
Le reti possono essere classificate in base alla loro dimensione geografica:
\begin{itemize}
    \item \textbf{PAN (Personal Area Network)}: connessioni di dispositivi molto vicini tipo all'interno di una stanza oppure la chiave di un automobile con apertura al tocco.
    \item \textbf{LAN (Local Area Network)}: connessione tra dipositivi all'interno di un edificio, è la tipologia più diffusa perchè offre ottimi compromessi
    \item \textbf{MAN (Metropolitan Area Network)}: sono le connessioni urbane, vicolate dai tronchi di fibra otica soot il terreno
    \item \textbf{WAN (Wide Area Network)}: connessioni molto ampie come reti nazionali e internazionali.
    \item \textbf{Internet}: è la rete globale che è formata da tante tipologie di reti differenti ma tutte conformi ai protocolli usati per la trasmissione e ricezione di informazioni.
\end{itemize}

Per comunicare sul mezzo di rete (cavo fisico/onde radio) i calcolatori devono interfacciarsi con la \textbf{scheda di rete}, essa è personale e presenta il proprio \textbf{MAC address} che è 
un codice di 48 bit unico per ogni scheda.
\paragraph{Evoluzione delle reti}
La nascita di internet è dovuta allo sviluppo globale delle reti, volendo tracciare una data di inizio essa ricadrebbe nel 1969 quando un gruppo di ricercatori universitari collegarono tra loro 4 
computer di 4 università americane; da quel momento molte entità, provate e pubbliche, hanno contribuito allo sviluppo e distribuzione delle reti.
Nel 2003 il mondo realizzò che internet poteva essere un'enorme vetrina commerciale veisibile da tutti gli utenti connessi, questi sono gli anni delle \textit{dot com} e si contavano circa 172 milioni
 di calcolatori connessi. Ad oggi il numero di dispositivi connessi risultano essere oltre i 4 miliardi, 25 se si considerano nche quelli dell'internet of things.
Il denaro usato per la realizzazione e manutenzione delle infrastrutture di rete viene principalmente dalle grandi multinazionali che lo recuperano facendo pagare l'acesso alla rete globale (\textit{last mile}) 
al singolo utente attraverso tariffe mensili.
\paragraph{Le prestazioni di rete}
Un tema estremamente importante sono le prestazioni con cui le reti di calcolatori riescono a trasmettere e ricevere, due sono le miusre importanti a riguardo:
\begin{itemize}
    \item \textbf{Capacità di trasmissione}: il numero di bit o byte trasmessi o ricevuti in un secondo, esso dipende principalmente dall'infrastruttura di rete.
    \item \textbf{Ritardo nel collegamento}: il tempo richiesto ai dati per transitre dal mittente al destinatario, esso è quindi la somma tra \textit{tempo di trasmissione} (durata del collegamenti) e \textit{ritardo intermedio}.
                                            esso è influenzato sia dalla distanza tra mittente e destinatario che dalla gestione intermedia della trasmissione (protocolli).
\end{itemize}

Un altro parametro molto importante è il \textbf{Jitter del ritardo di rete} ovvero la differenza, rispetto alla media, di ritardo tra i vari pacchetti. Questo parametro è molto importante perchè in base
all'utilizzo che si fa della rete si può decidere se prediligere un Jitter basso: se stessimo giocando ad un videogioco online l'interesse principale sarebbe avere il ritardo medio del pacchetto il 
più basso possibile, senza particolare preoccupazione se i pacchetti hanno una distanza molto alta dalla media. Viceversa, se stessimo guardando un contenuto in streaming, al fine di mantenere il 
\textit{buffer} più pieno possibile un Jitter basso è essenziale: se so che i pachetti mi arrivano molto vicini alla media del ritardo posso mantenere un buffer più ristrestto.
Ultimo parametro molto importante è il \textbf{RTT (Round Trip Home) o Ping} ovvero il tempo che un pacchetto ci mette a fare andata e ritorno.

\paragraph{Componenti per la connessione di rete}
PEr connettersi ad una rete un calcolatore ha bisogno di un numero minimo di componenti, tra cui:
\begin{itemize}
    \item \textbf{Scheda di rete}: distpositivo hardware per la codifica, trasmissione, ricezione e decodifica delle informazioni. E' dotata di tutte le funzionalità di cui il PC può usufruire come
    firmware, API garantiti dal produttore.
    \item \textbf{Mezzo di trasmissione}: supporto fisico alla trasmissione dei segnali, realizza l'infrastruttura di rete.
    \item \textbf{Connettore di rete}: è un connettore standard (RJ45) per collegare il dispositivo al mezzo di trasmissione di rete.
    \item \textbf{Protocollo di rete}: software del calcolatore atto a garantire la compatibilità tra tutti i dispositivi connessi alla rete.
\end{itemize}

\paragraph{Collegamenti e infrastrutture di rete}
Per attuare il collegamento di rete è necessario un mezzo di trasmissione comune tra tutti i dispositivi della retec che supporti la trasmissione fisica dei segnali secondo una codifica stabilita a priori.
L'\textbf{infrastruttura di rete} è la struttura dei collegamenti tra i calcolatori della rete, essa può essere:
\begin{itemize}
    \item \textit{Punto a punto}: è la struttura più semplice, l'unica difficoltà è che non possono parlare entrambi i calcolatori nello stesso momento.
    \item \textit{Completamente connese}: ogni calcolatore è connesso in modo diretto con ogni altro, oltre alla comunicazione diretta è possibile una comunicazione in più step. Il grande problema è che non 
    è per nulla scalabile.
    \item \textit{Parzialmente connesse}: ogni calcolatore è connesso solo con altri due calcolatori tipo corda, si usano poche risorse ed è sclabile ma la possibilità di trasmissione dipende dagli altri 
    calcolatori (mutua esclusione) e in caso di problemi diventa una partizione di rete.
    \item \textit{Partizione di rete}: gruppi di calcolatori isolati tra di loro, da evitare assolutamente.
\end{itemize}

\paragraph{Topologie di rete}
Le strutture di rete possono assumere vari schemi di connessione dei dispositivi, queste disposizioni sono chiamate topologie:
\begin{itemize}
    \item \textbf{Anello}: Il numero di collegamenti neccessario e di 1 superiore al minimo, questo rende però molto meno complicata la parte di spedizione perchè una volta che mi ritorna il messaggio mandato
    so che esso è stato ricevuto. Più precisamente per la trasmissione dei dati si usa un \textbf{token} ovvero una serie di bit che quando ricevuta rende possibile al calcolatore trasmettere i dati, inoltre chiamato Q il
    tempo che un calcolatore può tenere il token prima di ritrasmetterlo si può calcolare dopo quanto il token tornerà o se è stato perso. Questa topologia può supportare una rotture senza creare delle partizioni
    di rete ma il bottleneck è proprio il token: se esso viene corrotto durante la trasmissione è perso e non si può far generare, indipendentemente, ai calcolatori un nuovo token; inoltre è facilmente implementabile
    un hack al protocollo che permette di tenere per più tempo il token.
    \item \textbf{Stella}: Ha il minimo numero di collegamenti necessari ma alla prima rottura si forma almeno una partizione di rete, nel caso si rompa il calcolatore centrale però ognuno rimane isolato.
    Per i singoli host sembra un trasmissione punto a punto però, a livello fisico, può cambiare molto dal tipo di device installato:
    \begin{itemize}
        \item \textit{HUB}: manda elettricamente il segnale a tutti i dispositivi collegati alla rete.
        \item \textit{Switch}: riceve i pacchetti, sale di un livello del protocollo, e capisce dove mandarli e li trasmette solo alla porta giusta; può supportare il parallelismo e in generale è un
        dispositivo un pò più nobile dell'HUB.
    \end{itemize}
    Questo tipo di topologia, ad oggi, è una di quelle più utilizzate.
    \item \textbf{Bus}: è stata la prima topologia implementata, ha delle criticità tipo il fatto che se si rompe il bus si interrompono tutte le connessioni, i segnali perdono di intensità sulle  lunghe 
    distanze e due dispositivi non possono comunicare nello stesso momento (\textbf{collisione}). Per evitare collisioni inizialmente si ascoltava la dorsale: se era libera si trasmetteva se inveece era 
    occupata si aspettava un tempo random e si riprovava ma anche questo non proteggeva al 100\% dalle collisioni per cui servirebbe un token ma a quel punto risulterebbe un anello fatto male.
    \item \textbf{Albero}: \textit{topologia minimale} (se si rompe un filo qualcuno resta escluso), non ha vantaggi rispetto alle altre se non che garantisce che la distanza media tra calcolatori sia minima.
\end{itemize}

\paragraph{Mezzo fisico di trasmissione}
Il mezzo utilizzato per la trasmissione delle informazioni a livello fisico può essere di tre tipi:
\begin{enumerate}
    \item \textit{Cavetti o fili metallici}: trasmettono segnali elettrici, sono molto usati perchè affidabili e con buon compromesso tra costo e prestazioni.
    \item \textit{Fibre ottiche}: trasmettono segnali luminosi tramite la riflessione totale, hanno ottime prestazioni ma anche un costo considerevole.
    \item \textit{Wireless}: trasmissione attraverso le onde elettromagnetiche, permettono mobilità ma il collegamento non è sempre affidabile, capacità e distanza sono ancora abbastanza limitate.
\end{enumerate}

\paragraph{Scheda di rete}
E' il componente del calcolatore che realizza la connessione fisica, attraverso un connettore, alla rete e il suo compito è quello di trasmettere e ricevere i dati.
Questo componente ha un codice identificativo chiamato \textbf{MAC (medium access control)} chè è un codice di 48 bit scritti in esdecimale, si riconosce un indirizzo MAC perchè è scritto sempre come 6 coppie di cifre decimali
separate dai ":". La scheda di rete ha uno stack che gestisce tutte le trasmissioni, le quali possono anche avere destinatari diversi; generalmente distinguiamo due categorie di trasmissioni:
\begin{itemize}
    \item \textbf{Unicast}: trasmissione verso un solo destinatario.
    \item \textbf{Broadcast}: trasmissione verso tutti i destinatari, anche se è poco probabile che tutti riescano a ricevere il messaggio per cui il MAC address standard per una trasmissione broadcast è: FF:FF:FF:FF:FF:FF
\end{itemize}

\paragraph{Canali di comunicazione}
Un canale di comunicazione è una visione astratta del mezzo di trasmissione, distinguiamo il \textbf{canale punto a punto} da quello \textbf{ad accesso multiplo}.
Il canale punto a punto è un canale riservto alla comunicazione tra due soli host mentre quello ad accesso multiplo permette la connessione tra più host sullo stesso canale; questo però richiede 
arbitraggio di trasmissione per evitare collisione e l'inidirizzamento, per fare in modo che ogni pacchetto arrivi al giusto destinatario. Questo obbliga la definizione di un \textit{protocollo}.

\paragraph{Modalità di trasferimento}
Una serie di canali in cascata può realizzare un servizio di comunicazione in due modi diversi:
\begin{enumerate}
    \item \textbf{Commutazione di circuito}: viene riservato un circuito di canali per la singola comunicazione punto a punto,  il ritardo di connessione è basso ma si utilizza poco anche la risorsa di rete.
    Non è difficile da implementare, l'unica difficoltà è quella di trovare quale sia il circuito vero e proprio da riservare.
    \item \textbf{Commutazione di pacchetto}: i dati da trasmettere vengono divisi in pacchetti indipendenti e spediti sul canale separatamente. Ogni pacchetto conserva diverse informazioni tra cui 
    l'indirizzo del destinatario e i canali sono condivisi. Le risorse di rete vengono sfruttate al meglio in caso di canale condivisione ma ne 
    perde la velocità di connessione. Internet si basa sulla commutazione di pacchetto.
\end{enumerate}

\paragraph{Servizi orientati alla connessione}
Nelle reti a commutazione di pacchetto il servizio di trasmissione ha un ruolo determinato dei protocolli di rete utilizzati: ogni nodo intermedio immagazzina e rispedisce i pacchetti di dati. Ci sono due tipi di servizi:
\begin{itemize}
    \item \textit{Orientati alla connessione}: garantiscono la consegna ordinata dei pacchetti ricevuti e ritrasmettono quelli che sono stati eventualmente persi.
    \item \textit{Non orientati alla connessione}: i pacchetti possono seguire strade diverse, arrivare in ordine errato o non arrivare mai. Internet è fatto così perchè il protocollo TCP riesce a rendere
    la rete affidabile gestendo le problematiche.
\end{itemize}

Il collo di bottiglia è il router, soprattutto se tutti i pacchetti ci devono passare, per evitare che si perdano i pacchetti il mittente regola la velocità di trasmissione in base a 
quante conferme di ricezione gli arrivano (\textbf{windowing}).


\paragraph{I protocolli di rete}
Per protocollo si intende l'insieme di regole e procedure di gestione dei processi di comunicazione (\textit{come devo fare le cose}) e l'insimee di regole e formati per lo scambio di messaggi non ambigui
(\textit{come sono disposti i dati nei pacchetti}). A livello di architettura i protocolli di rete sono divisi in livelli, ogni livello risolve un problema di comunicazione e rimane estraneo agli altri 
comunicando attraverso un interfaccia per richieste e servizi. La compartimentazione dei protocolli è molto utile per renderli felssibili ai casi d'uso: si possono aggiungere/rimuovere livelli.
Il riferimento standard per definire l'architettura dei protocolli di rete è lo \textbf{Standard ISO/OSI RM} che organizza i protocolli in sette livelli:
\begin{etaremune}
    \item \textbf{Livello Applicazione}: è il livello che mette a disposizione delle singole applicazioni utente i metodi per connettersi alla rete.
    \item \textbf{Livello Presentazione}: serviva a tradurre i messaggi quando non erano identici i metodi di lettura bit (big/little endian), oggi in disuso.
    \item \textbf{Livello Sessione}: all'inizio della connessione identificava il client in modo univoco in modo tale che se la connessione veniva perduta non bsognava ritrasmettere da capo, oggi è
    in disuso perchè alle applicazioni conviene profilare gli utenti.
    \item \textbf{Livello trasporto}: è il livello fondamentale oggi realizzato da TCP/UDP, rende la connessione affidabile e controlla la congestione del traffico
    \item \textbf{Livello rete}: è il livello massimo a cui lavorano anche i router, serve per l'instradamento ed è realizzato grazie alle gerarchie \textbf{IP} (internet protocol).
    \item \textbf{Livello MAC/LCC}: è il livello che controlla l'accesso al medium: si occuppa della trasmissione evitando le collisioni.
    \item \textbf{Livello fisico}: generazione dei segnali.
\end{etaremune}

\paragraph{architettura dei protocolli internet}
L'architettura dei livelli di protocollo di internet utilizza solo 5 dei 7 livelli dello standard, in trasmissione ogni livello riceve dai dal livello superiore e gli incapsula con ulteriori dati che servono
a istruire il livello corrispondente del ricevitore. Alla ricezione ogni livello decapsula i dati leggendo le proprie istruzioni e li trasmette ai livelli superiori.

\begin{etaremune}
    \item \textbf{Livello Applicazione}: è il livello che mette a disposizione delle singole applicazioni utente i metodi per connettersi alla rete.
    \item \textbf{Livello trasporto}: imbusta i dati aggiungendo informazioni utili all'ordinamento e controllo della velocità d'invio. Il pacchetto è chiamato \textbf{segmento}.
    \item \textbf{Livello rete}: frammenta i dati in pacchetti e decide il cammino sul quale inviare il pacchetto.
    \item \textbf{Livello MAC/LCC}: esegue la consegna dei dati ai dispositivi di una rete locale capendo se quando è il momento più opportuno per evitare collisioni, il pacchetto è chiamato \textbf{frame}.
    \item \textbf{Livello fisico}: generazione dei segnali.
\end{etaremune}
Durante l'incapsulamento il livello trasporto scrive l'indirizzo IP del mittente e quello del destinatario, successivamente il livello MAC scrive l'indirizzo MAC del trsmettitore e del diretto ricevente.
In questo modo nel livello MAC cambiano sempre gli indirizzi perchè i vari router scrivono i loro inditizzi; questo porta alla presenza di 2 \textbf{Acknoledgment} uno fatto con i MAC address da ruouter a router, 
utili per evitare di dover rimandare i pacchetti da capo nel caso un router sia congestionato, e uno fatto con gli indirizzi IP nel caso il pacchetto si sia proprio perso e debba essere ritrasmesso dall'inizio.
inoltre ogni livello che incapsula insierisce un \textbf{header} con tutte le informazioni più rilevanti all'inizio e un \textbf{FCS (frame control system)} in coda.

\paragraph{Lielli e integrazioni delle reti}
In questo paragrafo vengono presentate le problematiche di ogni livello e come viene vista la rete da ogni livello:
\begin{etaremune}
    \item \textbf{Livello Applicazione}: la rete esiste come viene intesa da noi ed è utilizzabile dalle applicazioni utente
    \item \textbf{Livello trasporto}: la rete è una collezione di reti organizzate gerarchicamente, ha regole per la spedizione affidabile dei pacchetti e controlla la congestione di rete spingendo 
    sempre al massimo la portata della rete per poi rallentare.
    \item \textbf{Livello rete}: la rete è una collezione di reti organizzate gerarchicamente, ci sono regole per indirizzi di rete che nascondono i dettagli locali e i rputer sono i rappresentanti della ret locale.
    \item \textbf{Livello MAC/LCC}: la rete è locale, il livello può interagire con i mezzi trasmissivi con regole per indirizzi dei dispositivi, tempi di accesso e gestione errori
    \item \textbf{Livello fisico}: la rete è il mezzo di trasmissione, ha regole per la codifica e trasmissione dei dati che vede come segnali.
\end{etaremune}







\end{document}